\section{Introduction}
\label{sec:intro}

\IEEEPARstart{S}{elf}-supervised representation learning has transformed
language and vision, and is now reshaping electroencephalography (EEG)
analysis. Large EEG foundation models---LaBraM~\cite{labram2024},
EEGPT~\cite{eegpt2024}, NeuroLM~\cite{neurolm2025},
BIOT~\cite{biot2023}---pretrain on thousands of hours of EEG with 5--10\,M$+$
parameters on multi-GPU clusters, and transfer via fine-tuning or linear
probing to clinical and brain--computer interface (BCI) tasks. These
methods establish that scale helps, but their compute and parameter budgets
are incompatible with on-device deployment in wearable EEG hardware (Muse,
OpenBCI, Neurosity Crown, Smartbuds-class earbud devices) where latency,
energy, and memory budgets are measured in milliseconds, milliwatts, and
megabytes respectively.

\emph{When does compact SSL pretraining actually help on EEG, and how should
practitioners spend a limited compute budget?} The contemporary literature
provides partial answers. LeJEPA~\cite{lejepa2025} introduced the Sketched
Isotropic Gaussian Regularization (SIGReg) anti-collapse objective and a
small JEPA-style world model demonstrating $\sim$15\,M parameters trained on
a single GPU could be competitive at world-modelling. Concurrent work
Laya~\cite{laya2026} applied LeJEPA to EEG at foundation scale (29{,}000
hours of EEG, 20{,}940 subjects, two L40S GPUs), establishing that latent
prediction outperforms reconstruction-based pretraining on clinical EEG
benchmarks under \emph{cross-dataset evaluation}: EEGMMIDB and BCI Competition~IV
are held out from Laya's pretraining and pooled into multi-dataset downstream
benchmarks, on which no SSL foundation model exceeds 0.51 balanced accuracy
on left-versus-right motor imagery~\cite[Tab.~1]{laya2026}. The complementary
\emph{in-distribution single-corpus regime}---relevant to many applied EEG
settings where a research lab pretrains and evaluates on a focused corpus
realistic to its experimental design---is not addressed. Neither paper
characterizes the small-scale scaling regime, the SIGReg-weight optimum, or
the cross-dataset channel-transfer behaviour that practitioners face when
deploying a single pretrained encoder across heterogeneous downstream
recording rigs.

This paper fills that gap. We present \textbf{EEG-LeJEPA}, a 2.86\,M-parameter
encoder-plus-causal-predictor architecture trained with
SIGReg~\cite{lejepa2025} on 50 subjects of EEGMMIDB motor-imagery data---a
deliberately compact-scale recipe runnable on a single consumer GPU in 10
minutes. Across six (data, compute) operating points we map the scaling law,
we explicitly characterize the SIGReg-weight optimum, and we compare two
strategies for cross-dataset transfer (channel padding vs.\ channel-matched
pretraining). We compare against (i) a randomly-initialized identical
architecture as the no-pretraining control and (ii) the same architecture
trained end-to-end with supervised cross-entropy as the no-SSL control.

\subsection*{Contributions}

\begin{enumerate}
\item \textbf{A leakage-clean evaluation protocol for cross-subject EEG
self-supervised learning.} Under a strictly disjoint pretraining/evaluation
subject split, predictor-derived feature sources lose their apparent
advantage and the simple encoder mean-pool is the strongest and only
consistently significant source on both EEGMMIDB binary tasks
(left-vs-right MI $+10.9$\,pp; rest-vs-activity $+11.7$\,pp; both $q<0.002$
after Benjamini--Hochberg correction). The protocol is directly reusable
for any EEG-SSL study.

\item \textbf{The SIGReg projection-count rule.} Holding architecture,
data, and step count fixed and varying \emph{only} the Cram\'er--Wold
projection count $M$, a 2.4$\times$ larger encoder (256-d embedding)
collapses on left-vs-right MI from 0.650 ($M{=}1024$) to 0.567
($M{=}256$) and drops \emph{below} the 192-d base model, reproducing a
spurious ``larger-model regresses'' artifact. $M$ must scale with
embedding width---a previously unreported practical pitfall when applying
SIGReg-style isotropy regularizers across model widths.

\item \textbf{Characterization of the compact, single-corpus, single-GPU
regime.} Across six (data, compute) operating points (3--89 subjects,
200--30k training steps) under the leakage-free split, best-source AUC
rises and plateaus near 0.76; at the plateau, a 2.4$\times$ larger model
is modestly but consistently better at matched data ($+1.3$ to
$+2.0$\,pp on encoder mean-pool, seed-stable across three seeds), so the
compact model is a deliberate accuracy/size trade-off rather than a
capacity-optimal point. End-to-end pretraining$+$probe runs in under
15\,minutes on a single consumer GPU at total cost below \$1, with the
EEGMMIDB headline tasks reaching $68.6\,\%/0.751$\,AUC
(rest-vs-activity) and $62.1\,\%/0.690$\,AUC (left-vs-right MI).

\item \textbf{A counterintuitive cross-dataset transfer recipe.}
Channel-padded transfer from the full-channel pretrained encoder
\emph{outperforms} channel-matched re-pretraining ($+4.9$ vs.\
$+1.5$\,pp on encoder mean-pool, EEGMMIDB$\to$BCI-IV-2a). For
practitioners deploying a single pretrained encoder across heterogeneous
montages, this argues for pretraining at the richest available montage
and zero-padding at inference rather than re-pretraining per montage.

\item \textbf{A SIGReg-weight ablation with empirical anti-collapse
validation.} Across $\lambda \in \{0, 0.1, 0.3, 1.0, 3.0, 10.0\}$ the
downstream-accuracy curve is sharply U-shaped with an optimum at
$\lambda{=}1.0$ and \emph{complete representational collapse} at
$\lambda{=}0$ (probe accuracy regresses to exactly chance), complementing
the theoretical analysis of~\cite{lejepa2025} with a direct empirical
demonstration.

\item \textbf{A direct comparison to supervised-from-scratch training}
of the identical architecture: the SSL$+$linear-probe pipeline ties
supervised training on left-vs-right MI ($-0.2$\,pp, n.s.) and
significantly exceeds it on rest-vs-activity ($+8.4$\,pp, $p<0.001$),
while a single frozen encoder serves both downstream tasks at sub-second
per-task probing cost.

\item \textbf{A complete reproducibility package.} 39 unit tests across
six test modules; trained checkpoints (six $\lambda$-ablation checkpoints
plus the scaling, capacity multi-seed, $M$-isolation, and channel-transfer
checkpoints); per-fold CSV outputs for every reported table; the
dependency-free significance scripts; and end-to-end training$+$probe
runtime under 15 minutes on a single consumer GPU. Code, configurations,
and checkpoints released at \url{https://github.com/lifeinpassion/eeg-lejepa}.
\end{enumerate}

The remainder of the paper is organized as follows.
\cref{sec:related} surveys EEG foundation models and SSL anti-collapse
mechanisms with explicit relation to concurrent work
\cite{laya2026}. \cref{sec:method} presents the EEG-LeJEPA architecture
and SIGReg objective. \cref{sec:experiments} details the empirical setup
and downstream evaluation protocols. \cref{sec:results} presents the
scaling, ablation and transfer results. \cref{sec:discussion} situates
our findings against published EEG foundation models and outlines
limitations and future work. \cref{sec:conclusion} concludes.
